\section{Physical response geometry and limits of elimination}\label{app:geometry}

\subsection{A telescoping certificate on a path}\label{s6:telescoping}
For $0<\epsilon<1/2$, let
\begin{equation}
 P_\epsilon=\{x\in\R^n:0\le x_1\le1,\quad
 \epsilon x_{j-1}\le x_j\le1-\epsilon x_{j-1}\ (2\le j\le n)\},
 \qquad x_0=0.\label{s6:km}
\end{equation}
Every coordinate lies in $[0,1]$. The two inequalities in each pair
cannot both be tight. Thus a vertex has one tight inequality in each
pair; conversely any such choice gives a nonsingular triangular system
and a feasible vertex. Its coordinates are
$x_j(u)=u_j+(1-2u_j)\epsilon x_{j-1}(u)$ for $u\in\{0,1\}^n$.
All $2^n$ terminal values are distinct: the last bit chooses the disjoint
intervals $[0,\epsilon]$ and $[1-\epsilon,1]$, and within either interval
its affine recursion can be inverted, successively recovering all bits.

\begin{lemma}[A vertex certificate]\label{s6:certificate}
Put
\[
 d_j=(1-\epsilon)\epsilon^{2(n-j)-1}\quad(j<n),\qquad
 L(x)=x_n-\sum_{j<n}d_jx_j,\qquad F(x)=L(x)-x_n^2.
\]
Then
\begin{equation}
 F(x)=\sum_{j=1}^n\epsilon^{2(n-j)}
 (x_j-\epsilon x_{j-1})(1-\epsilon x_{j-1}-x_j)\ge0
 \quad(x\in P_\epsilon),\label{s6:telescoping-identity}
\end{equation}
with equality precisely at the vertices. Every vertex $v$ is uniquely
exposed by $\sum_{j<n}d_jx_j+(2v_n-1)x_n$.
\end{lemma}
\begin{proof}
The $j$th product expands to
$x_j-x_j^2-\epsilon x_{j-1}+\epsilon^2x_{j-1}^2$.
The weighted squares telescope, leaving $-x_n^2$, and the remaining
coefficients give $L$. All factors and weights are nonnegative, so equality
requires a tight endpoint in every pair, exactly the vertex condition.
For $t=v_n$, throughout $P_\epsilon$,
\[
 2tx_n-L(x)\le 2tx_n-x_n^2=t^2-(x_n-t)^2\le t^2.
\]
Equality requires $F=0$ and $x_n=t$, which select $v$ uniquely.
\end{proof}

Exponential two-dimensional shadows of this Klee--Minty cube are established
by G\"artner, Helbling, Ota and Takahashi
\cite[Section~4]{s6:gartner2013}. The lemma supplies a convenient
self-contained certificate for the physical constructions below; it is
not a claim that the exponential-shadow phenomenon is new.
For comparison, their direction at $\epsilon=1/4$ uses
$c_j=\epsilon^{3(n-j)}$ for $j<n$, $c_n=0$, and exposes every vertex
by $c^Tx+\lambda x_n$ with $|\lambda|<1/15$.
The direction $d$ in the lemma instead uses the interval $[-1,1]$.
Both directions have polynomial rational encoding length.

\subsection{Physical paths and one varying product price}\label{s6:physical-paths}
Fix $n\ge2$, $\epsilon=1/4$, and $s_j=4^{j-n}$. Source $j$ has
quality $C_j=n-j$ and exact supply $s_j$. Its only arcs go to products
$j-1$ and $j$, with flows $s_j-t_j$ and $t_j$, each bounded above by
$s_j$. Product~0 has capacity $s_1$ and redundant upper quality $C_1$;
product~$n$ has capacity one and redundant upper quality $C_n=0$.
For $2\le j\le n$, product $j-1$ has capacity $s_j$ and upper quality
$(C_{j-1}+C_j)/2$.

\begin{proposition}[Exact path realization]\label{s6:path-realization}
This pool-free physical network is affinely bijective to $P_{1/4}$ by
$t_j=s_jx_j$. Its bipartite bypass graph is one path, and every source
and product has degree at most two. All upper flow bounds are at most
one; scaling all qualities by $1/(n-1)$ puts them in $[0,1]$.
\end{proposition}
\begin{proof}
The internal product receives $t_{j-1}+s_j-t_j$. Its capacity is
$t_j\ge t_{j-1}$. Subtracting its midpoint quality from the two
incident source qualities gives opposite coefficients $1/2$ and
$-1/2$, so its homogeneous quality row is
$t_{j-1}\le s_j-t_j$. Thus its two rows are exactly
\begin{equation}
 t_{j-1}\le t_j\le s_j-t_{j-1}.\label{s6:physical-strip}
\end{equation}
After dividing by $s_j$, the predecessor coefficient is $1/4$.
The endpoint and arc bounds give the other bounds, and the reverse
substitution satisfies every physical row. All other flows are uniquely
$s_j-t_j$; hence the correspondence is bijective. Scaling quality does
not change any homogeneous quality inequality.
\end{proof}

\begin{theorem}[Exponential physical price response]\label{s6:price-response}
There is a family of rational blending paths with $O(n)$ nodes, one
upper-bounded quality, all upper flow capacities at most one, zero
production costs, and positive product revenues, for which varying just
the final product's unit revenue exposes $2^n$ distinct unique optimal
flows. The family can have only two distinct input qualities and no
positive lower flow bounds. Its optimal-value function has $2^n$
open affine pieces. Every fixed-price problem is an LP.
\end{theorem}
\begin{proof}
First keep the exact source contracts in \cref{s6:path-realization}.
Use the direction $d$ from \cref{s6:certificate}; its coefficient in
physical rightward flows is $w_j=d_j/s_j=3s_j$ for $j<n$.
Put $w_n=\lambda\in[-1,1]$ and define output revenues
\[
 R_0=1,\qquad R_j=1+\sum_{h=1}^j w_h\quad(1\le j\le n).
\]
Since $\sum_{j<n}3s_j=1-s_1<1$, all revenues lie in $[0,3]$ and only
$R_n$ depends on $\lambda$. With zero costs their total revenue is
\begin{equation}
 \sum_j[R_{j-1}(s_j-t_j)+R_jt_j]
 =C_0+\sum_{j<n}d_jx_j+\lambda x_n,
 \quad C_0=\sum_jR_{j-1}s_j.\label{s6:price-identity}
\end{equation}
The constant $C_0$ is independent of $\lambda$. The lemma exposes every
vertex, including $x_n=0,1$ at $\lambda=-1,1$. Strict optimality among a
finite vertex set persists on a neighborhood of each exposing parameter,
so each has a nonempty one-sided interval inside $[-1,1]$, and an open
interval in its interior. Distinct terminal values give distinct slopes;
there are exactly $2^n$ affine regimes. One can instead use the published
direction $c$ and $w_j=c_j/s_j=16^{j-n}$, obtaining revenues in $[0,2]$
and the smaller price interval $[-1/15,1/15]$.

To use only two input qualities, replace the descending quality path by
the following relay path. Start with a quality-one source of supply
$s_1$. For each $j=2,\ldots,n$, attach a quality-zero source of supply
$s_j$ through a product with capacity $s_j$ and upper quality $1/2$.
If $j<n$, append a quality-one source of the same supply $s_j$ through
a product with exact demand $s_j$ and redundant upper quality one.
A descending step gives \eqref{s6:physical-strip}. At a reset product,
\[
 t_j+(s_j-\widetilde t_j)=s_j
 \quad\Longleftrightarrow\quad\widetilde t_j=t_j.
\]
Thus the quality-one duplicate copies the signal before the next
quality-zero source. End products have redundant upper quality one
and capacities $s_1,1$. There are $m=2n-2$ sources with supply sequence
\[
 S_1,\ldots,S_m=(s_1,s_2,s_2,s_3,s_3,\ldots,s_{n-1},s_{n-1},s_n).
\]
All arcs of source $i$ have upper bound $S_i$.
The original free rightward flows and their equal duplicates determine
all physical flows, giving the same affine bijection to $P_{1/4}$.
Assign coefficient $3s_j$ to the free-coordinate position $j<n$,
zero to reset duplicates, and $\lambda$ to the last position.
With these coefficients in path order denoted $g_i$, take
$\rho_0=1$, $\rho_i=1+\sum_{h\le i}g_h$.
The same telescoping revenue computation gives
$\sum_i\rho_{i-1}S_i+\sum_{j<n}d_jx_j+\lambda x_n$.
Only the final revenue varies and every base revenue is in $[0,3]$.

It remains to remove all source and reset-product lower contracts.
Let $Q$ be the nonempty exact-contract polytope in its $N=2m$ arc
variables $f$. Write these contracts $Ef=e$. Retain their upper bounds
and drop their lower bounds, so the deficit vector
$\Delta=e-Ef$ is nonnegative; put $\delta=\mathbf1^T\Delta$.
All coefficient rows, after clearing the midpoint factor two, have
entries in $\{-1,0,1\}$. By \cref{s3:hoffman}, every relaxed flow has
an exact-contract repair $f'\in Q$ with
$\|f'-f\|_1\le H\delta$, where $H=N^N$.
The base revenue vector has infinity norm at most three, uniformly in
$\lambda$. Choose $M=3H+1$, add $M$ to every output revenue, and add
another $M$ to every reset-output revenue. Conservation makes the first
bonus reward all source throughputs; the second rewards reset throughputs.
Writing $E_0=\mathbf1^Te$, the new objective $B(f)+M(E_0-\delta)$ obeys
\[
 B(f)+M(E_0-\delta)
 \le B(f')+ME_0-(M-3H)\delta.
\]
Thus every optimizer satisfies every removed contract. The value function
changes by exactly $ME_0$ and every previously exposed optimizer remains
unique. Revenues are now positive, all production costs remain zero,
and $M$ has $O(n\log n)$ bits. No bound on revenue magnitude is claimed.
For the original direction $c$, the same proof uses $M=2H+1$.
The version with varying input qualities needs no reset contracts and
uses $N=2n$ and the identical source-throughput bonus argument.
\end{proof}

\begin{corollary}[Explicit elimination needs exponential output]\label{s6:response-size}
For either response in \cref{s6:price-response}, a piecewise-affine table
requires $2^n$ pieces. Any quantifier-free Boolean formula in price and
value alone describing its graph, or its epigraph on the price interval,
uses polynomials whose total degree is at least $2^n$.
The fixed-objective state message
$G(t)=\max\{\sum_{j<n}d_jx_j:x\in P_{1/4},\ x_n=t\}$ has
$2^n-1$ affine intervals.
\end{corollary}
\begin{proof}
Remove identically zero polynomials from a proposed formula. Where none
of its remaining polynomials vanishes, all signs, and hence its truth
value, are locally constant. Each open affine graph segment is a
boundary segment, so the product of the defining polynomials vanishes
on that segment. Restriction to its supporting line vanishes on an
interval and therefore identically; the line equation divides the
product. The $2^n$ distinct lines require that many linear factors.
This applies also after any upper clipping of the epigraph that leaves
all those segments below the clip.
For $G$, all projected vertex points $(x_n,\sum d_jx_j)$ are uniquely
exposed with coefficient $+1$ on the second coordinate. They lie on its
upper boundary, have distinct first coordinates, and consecutive points
cannot be collinear. Their first coordinates run from zero to one,
giving exactly $2^n-1$ intervals.
\end{proof}

These are limits on explicit descriptions. Backward elimination still
evaluates the support at a rational price with $O(n)$ arithmetic operations:
set $a_n=\lambda$ and $a_j=d_j-\epsilon|a_{j+1}|$ for $j=n-1,\ldots,1$;
the value is $\sum_j\max(0,a_j)$. Eliminating a variable with positive
current coefficient selects its upper endpoint, adds that coefficient
to the constant, and subtracts $\epsilon$ times its absolute value from
the predecessor coefficient; a negative coefficient selects its lower
endpoint and makes the same predecessor update without the constant.
All intermediate bit lengths are polynomial. The same recurrence works
with $c$ instead of $d$. Compact circuits, extended formulations, and
implicit algorithms are therefore fully compatible with the lower bound.
In particular, \cref{s4:path-projection} projects onto endpoint coordinates,
whereas $G$ retains a dense eliminated objective.

\begin{proposition}[A limit on individual coordinate ports]\label{s6:ports}
Without nonphysical global rows, a pool-free network with source and
product degrees at most two cannot represent every bounded rational
polytope by a coordinate projection onto selected individual arc flows,
even with unrestricted attribute count and two-sided bounds.
\end{proposition}
\begin{proof}
Every physical row involves at most two arc variables. Fourier--Motzkin
elimination preserves this property: a pair of rows of opposite signs
in the eliminated variable contains at most one other variable each.
Thus every coordinate projection has a finite two-variable-per-inequality
description; its size need not be polynomial. The simplex
$\{x\in\R_+^3:x_1+x_2+x_3\le1\}$ has no such description.
The point $(1/2,1/2,1/2)$ satisfies every valid inequality involving at
most two coordinates, because each of its coordinate pairs extends to
a simplex point by setting the third coordinate to zero. Yet that point
is outside the simplex. This does not restrict arbitrary linear images.
\end{proof}

\subsection{One pool: many optima and an exact LP hull}\label{s6:geometry}
Use the contracted path of \cref{s6:path-realization}, retaining its
unscaled qualities $C_j=n-j$, and put $t=t_n=x_n$. Modify its terminal
product and add the following interface:
\[
 A\to n,\quad A\to P,\quad B\to P,\quad B\to W,\quad
 P\to W,\quad P\to V,\quad Z\to V.
\]
Sources $A,B$ have qualities one and two and exact supplies one;
source $Z$ has quality zero and supply interval $[0,1]$. Every new arc
has capacity one. Pool $P$ has exact throughput one. Products $n,W$
have exact demand one and upper qualities one and two; product $V$ has
capacity two and upper quality one. There are no further positive
lower bounds. The original terminal source has quality zero.
\Cref{s6:interface-figure} shows the interface and its forced flows.
\begin{figure}[!htbp]
\centering
\begin{tikzpicture}[>=Stealth,scale=.90,
 source/.style={circle,draw,minimum size=7mm,inner sep=1pt},
 product/.style={rectangle,draw,minimum size=7mm,inner sep=2pt},
 pool/.style={circle,draw,thick,minimum size=9mm,inner sep=1pt},
 every edge/.style={draw,->},every node/.style={font=\small}]
\node[source] (e) at (-2,1) {$n$};
\node[product] (n) at (0,1) {$n$};
\node[source] (a) at (2,1) {$A$};
\node[pool] (p) at (4,1) {$P$};
\node[source] (b) at (4,-1) {$B$};
\node[product] (w) at (6,-1) {$W$};
\node[product] (v) at (6,1) {$V$};
\node[source] (z) at (6,2.7) {$Z$};
\path (e) edge node[above] {$t$} (n)
(a) edge node[above] {$1-t$} (n)
(a) edge node[above] {$t$} (p)
(b) edge node[left] {$1-t$} (p)
(b) edge node[below] {$t$} (w)
(p) edge node[below left] {$1-t$} (w)
(p) edge node[above] {$t$} (v)
(z) edge node[right] {$z$} (v);
\node[below=2mm] at (e.south) {$C_n=0$};
\node[below=2mm] at (a.south) {$C_A=1$};
\node[below=2mm] at (b.south) {$C_B=2$};
\node[above=2mm] at (z.north) {$C_Z=0$};
\node[above=4mm] at (p.north) {$q=2-t$};
\end{tikzpicture}
\caption{The vertex-forcing interface attached to the terminal source
and product of the path. Circles are sources, the thick circle is the
pool, and squares are products. Contracts force the displayed flows;
the quality row at $V$ becomes $z\ge t-t^2$. The triangle $B,P,W$ is
the only undirected cycle. The terminal source's other path arc is
omitted.}
\label{s6:interface-figure}
\end{figure}


\begin{theorem}[A physical vertex-forcing family]\label{s6:vertex-forcing}
The resulting one-pool network has a feasible set affinely equivalent,
in physical arc flows and active pool quality, to
\begin{equation}
 T=\{(x,z):x\in P_{1/4},\quad x_n-x_n^2\le z\le1\}.
 \label{s6:T}
\end{equation}
There are ordinary nonnegative source costs and product revenues bounded
by two for which the network has exactly $2^n$ isolated global maximizers.
Every source and product has degree at most two, the pool has two feeds
and two outlets, and the full undirected network is connected with one
cycle. All quality data can be scaled into $[0,1]$.
\end{theorem}
\begin{proof}
Demand at $n$ forces $f_{An}=1-t$; source $A$ then gives $f_{AP}=t$.
Pool throughput gives $f_{BP}=1-t$, whence source $B$ gives $f_{BW}=t$.
Demand at $W$ forces $f_{PW}=1-t$, leaving $f_{PV}=t$.
The pool is active with quality $q=2-t$. If $z=f_{ZV}$, the only
nonredundant new quality row is
\[
 (2-t)t\le t+z\quad\Longleftrightarrow\quad z\ge t-t^2.
\]
Every $t\in[0,1]$ has $0\le t-t^2\le1/4$, so all other arc,
source, throughput and quality bounds hold, and this describes every
physical extension. In particular all flows and $q$ are affine in
$(x,z)$, despite the curved feasible boundary.

Give path source $j$ cost $s_j$, sources $A,B$ cost one, and $Z$ cost
two. Give path product $j<n$ revenue $s_{j+1}$ and products $n,W,V$
revenue one. Subtracting one from every source cost and every product
revenue leaves profit unchanged by total mass conservation. The new
interface prices are zero except for unit cost on $Z$. On the path,
source $j$ has cost $s_j-1$, its left output has the same revenue,
and its right output has revenue $s_{j+1}-1$ when $j<n$ and zero
when $j=n$. Thus profit is exactly
\begin{equation}
 \sum_{j<n}(s_{j+1}-s_j)t_j-z
 =\sum_{j<n}3s_j^2x_j-z=\sum_{j<n}d_jx_j-z.
 \label{s6:geometry-profit}
\end{equation}
By \cref{s6:certificate}, this is at most $-F(x)\le0$; equality forces
a path vertex and the unique clean flow $z=t-t^2$. These are all
$2^n$ optimizers, so all are isolated. The degree assertions follow from
the displayed arcs. The original path and its attached trees introduce
no cycle; the sole cycle is $B-P-W-B$. Dividing every quality and
specification by $\max(n-1,2)$ preserves all flows and normalizes data.
\end{proof}

\begin{corollary}[Strict local optima with distinct values]\label{s6:local}
A polynomially encoded nonnegative price perturbation of this network
has at least $2^n$ strict local maxima with distinct values, exactly one
of the constructed points being globally optimal. Its prices are at
most three.
\end{corollary}
\begin{proof}
Increase both the revenue of $V$ and cost of $Z$ by $\delta=4^{-n}$.
Their net contribution is $\delta[(t+z)-z]=\delta t$.
At least clean flow, the objective is $H(x)=-F(x)+\delta x_n$.
At a vertex of terminal value $t$ it is $\delta t$. Every terminal
vertex value has denominator dividing $4^{n-1}$, and these values are
distinct. For neighboring vertices $v,v'$ with terminal values $t,t'$,
linearity of $L$ and $F(v)=F(v')=0$ give
\[
 F((1-a)v+av')=a(1-a)(t'-t)^2.
\]
The outward derivative of $H$ is
$-(t'-t)^2+\delta(t'-t)<0$, since $|t'-t|\ge4^{-(n-1)}>\delta$.
The tangent cone at a polytope vertex is generated by its incident edge
rays. Strict negativity on these finitely many rays bounds the linear
term above by $-\eta\|x-v\|$ for some $\eta>0$ on that cone; the
quadratic remainder is $O(\|x-v\|^2)$. Hence the vertex is a strict
local maximum. Extra clean flow decreases profit strictly, so the
property transfers to \eqref{s6:T} and then to the physical coordinates.
Finally every feasible profit is at most $\delta t\le\delta$, with
equality only at the unique vertex with $t=1$ and least clean flow.
No assertion excludes additional local maxima. The perturbation shrinks
with dimension and supplies no uniform tolerance or local-solver runtime
bound.
\end{proof}

\begin{corollary}[Exact integer dimension of the optimal set]\label{s6:integer-lift}
Let $S$ be the complete optimal physical-flow set for the unperturbed
objective in \cref{s6:vertex-forcing}, optionally retaining the pool quality.
The minimum number of integer coordinates in any exact convex lift of
$S$ is $n$, allowing arbitrary continuous auxiliaries and unbounded
integer variables.
\end{corollary}
\begin{proof}
For distinct optimal path vertices $v,v'$,
$F((v+v')/2)=(v_n-v_n')^2/4>0$, so their physical midpoint is not
optimal. Suppose a convex set with $k<n$ integer coordinates projects
exactly to $S$ when those coordinates are integral. Choose a lift of
each of its $2^n$ points. Two integer vectors have the same parity;
their midpoint is integral and belongs to the convex lifting set.
Its physical projection is the forbidden midpoint, a contradiction.
For the upper bound, attach to the optimizer indexed by $u$ the vector
$u\in\{0,1\}^n$ and take the convex hull of these finitely many lifted
points. Any integral vector in its final $n$ coordinates is binary;
a convex combination yielding it must use only points with that same
binary vector. It therefore selects exactly its original optimizer.
This upper construction may have exponentially many vertices.
The lower-bound argument is the established midpoint/parity principle
of Lubin, Vielma and Zadik \cite[Section~4.2]{s6:lubin2022}, applied to
this physical optimal set.
\end{proof}

\begin{theorem}[A compact exact hull and polynomial optimization]\label{s6:lp-hull}
The same physical feasible family admits an exact polynomial-size LP
convex hull. Every rational linear objective in physical flows and active
pool quality has a polynomial-time exact physical optimizer, with rational
coordinates of polynomial bit length.
\end{theorem}
\begin{proof}
Writing $d^Tx=\sum_{j<n}d_jx_j$, the exact hull in \eqref{s6:T} is
\begin{equation}
 \conv(T)=\{(x,z):x\in P_{1/4},\quad d^Tx\le z\le1\}.
 \label{s6:hull}
\end{equation}
The certificate gives $d^Tx\le x_n-x_n^2$, proving containment.
Conversely write $x$ as a convex combination of path vertices.
At every vertex $v$, $v_n-v_n^2=d^Tv$, so $(x,d^Tx)$ is the same
convex combination of points in $T$. Also $(x,1)\in T$. Their vertical
segment contains every point on the right of \eqref{s6:hull}.
A vertex of that polytope must lie on $z=d^Tx$ or $z=1$, since otherwise
it admits a vertical perturbation. On either affine boundary, a nonvertex
$x$ admits a nontrivial convex decomposition. Consequently every hull
vertex has $x$ a path vertex and is physically feasible. The two
boundaries do not meet, since $d^Tx\le x_n-x_n^2\le1/4$.
Solve the rational LP and return an optimal vertex, then apply the affine
physical map. An arbitrary nonvertex point of its optimal face need not
be physically feasible. Polynomial bit length follows from rational LP
vertex bounds. The hull does not exactly represent the original
nonconvex set or its finite optimal set, so it does not contradict
\cref{s6:integer-lift}.
\end{proof}

\begin{corollary}[Three qualities and general supply sequences]\label{s6:alphabet-geometry}
The conclusions of \cref{s6:vertex-forcing,s6:local,s6:integer-lift,s6:lp-hull}
hold with only three distinct input qualities. The unperturbed geometry,
integer dimension and LP hull also extend to any rational supplies with
$s_n=1$, $s_1>0$ and $s_{j+1}>2s_j$ in the descending-quality path.
\end{corollary}
\begin{proof}
For the first assertion attach the identical interface to the contracted
two-quality relay path in the proof of \cref{s6:price-response}. Its
last source has quality zero, so every interface calculation is unchanged.
Keep all its source, reset-output and interface contracts. At physical
source position $i$ give cost $S_i$, left endpoint output revenue $S_1$,
and output after position $i<m$ revenue $S_{i+1}$; terminal and interface
prices remain as before. The coefficient of its rightward flow is
$S_{i+1}-S_i$. This coefficient vanishes when the next source is the
same-supply reset duplicate. The coefficient $3s_j$ before the next
descending step occurs on that duplicate, whose rightward flow equals
$t_j$. At the first position it occurs directly on $t_1$. Thus the
sum is again \eqref{s6:geometry-profit}. All conclusions transfer through
the affine bijection, and the full quality alphabet is $\{0,1,2\}$.

For the second assertion keep \eqref{s6:physical-strip}, which remains
the physical capacity/quality pair for these supplies. Direct expansion gives
\[
 t_n-t_n^2-\sum_{j<n}(s_{j+1}-s_j)t_j
 =\sum_{j=1}^n(t_j-t_{j-1})(s_j-t_{j-1}-t_j),\qquad t_0=0.
\]
Every summand is nonnegative. The strict growth of supplies makes both
endpoint branches distinct and their terminal ranges disjoint, so there
are exactly $2^n$ endpoint vertices and distinct terminal values.
The interface and economics now use coefficients $s_{j+1}-s_j$.
The equality, midpoint and convex-hull proofs apply verbatim in these
physical coordinates. Encoding and algorithmic bounds include the given
supply bit lengths. The specific perturbation $4^{-n}$ is asserted only
for the geometric sequence.
\end{proof}

The linear exact-penalty argument in \cref{s6:price-response} does not
remove the contracts in the nonlinear interface. In fact deleting just
the lower pool-throughput bound changes its optimum: set
$x_1=\cdots=x_{n-2}=0$, $x_{n-1}=1$, $x_n=1/4$, send $1/4$ from
$A$ through the pool to $V$, send the rest of $A$ to $n$, send all of
$B$ to $W$, and set $z=f_{PW}=0$. All other bounds and contracts hold,
the pool quality is one, and profit is $d_{n-1}=3/16>0$.
Multiple local optima in pooling have earlier geometric examples;
Grothey and McKinnon \cite[Section~3]{s6:grothey2020} describe a
three-pool, two-quality example. The present family adds the explicit
asymptotic count, degree restrictions and LP hull; this comparison does
not establish exhaustive priority for those refinements.

\subsection{Why one unrestricted aggregate changes the abstract problem}
\label{s6:slab}
The preceding geometry alone cannot prove degree-two pooling hardness.
The following abstract reduction identifies an additional ingredient.
Rank-one concave quadratic hardness in general is established by
Pardalos and Vavasis \cite{s6:pardalos1991}; the point here is its local
path structure and the role of a dense slab.

\begin{proposition}[A path with one dense slab]\label{s6:slab-hardness}
For rational $0<\epsilon<1/2$, deciding whether
\[
 \min\{F(x):x\in P_\epsilon,\ B-1/4\le w^Tx\le B+1/4\}\le0
\]
is ordinarily $\NP$-complete, even on nonempty compact instances.
The Hessian of $F$ has one negative eigenvalue and rank one. Local
path coefficients can alternatively be fixed to $0,\pm1,\pm1/4$,
with additional singleton upper bounds and polynomially many variables.
\end{proposition}
\begin{proof}
For a SUBSET SUM instance of positive integer weights $w_i$ and integer
$0\le B\le W=\sum_iw_i$, set $\epsilon=1/(8W)$. A vertex with
endpoint bits $u$ satisfies $|x_i-u_i|\le\epsilon$.
If $w^Tu=B$, then $|w^Tx-B|\le1/8$, so that vertex satisfies the
slab and $F=0$. Conversely $F\le0$ forces a vertex by
\cref{s6:certificate}; its bits satisfy
$|w^Tu-B|\le1/4+1/8<1$, forcing equality of the integers.
The zero vertex has weighted sum zero. The all-one-bit vertex has
weighted sum at least $W-1/8$. If $B<W$, convex interpolation reaches
$B$; if $B=W$, the latter vertex already lies in the slab.
Thus the domain is nonempty. The coefficients in $F$, including powers
of $\epsilon$, have polynomial bit length. The endpoint bits are an
$\NP$ certificate: reconstruct the rational vertex and verify the slab.
The Hessian is $-2e_ne_n^T$.

To fix $\epsilon=1/4$, let $r$ be the least nonnegative integer with
$4^{-(r+1)}\le1/(8W)$. Place free coordinates at positions
$1,r+2,2r+3,\ldots$ on a path of dimension $N=n+(n-1)r$.
Impose $x_j\le1/2$ on each padding coordinate and use the full
$N$-coordinate certificate $F$. At a zero of $F$, padding positions
must take their lower branch: the upper branch is at least $3/4$.
Successive free coordinates are separated by $r$ multiplications by
$1/4$, so each differs from its endpoint bit by at most $4^{-(r+1)}$.
Every free-bit vector extends uniquely to a zero of $F$. Put the slab
weights only on free positions; the same rounding, nonemptiness and
certificate arguments apply. Since $r=O(\log W)$ this is polynomial.
The dense weights and objective coefficients retain binary numerical
data, so this gives no strong-hardness or inverse-polynomial gap claim.
\end{proof}

Fixing the scalar nonlinear coordinate $q=x_n$ in this problem leaves
a bounded LP. The local inequalities form a path, but the slab has two
bounds on a dense aggregate, and $L$ is another dense linear form.
Thus one nonlinear parameter and a fixed number of aggregate rows do
not justify replacing the independent blocks of \cref{s4:fixed-core}
by arbitrarily long coupled paths. No physical degree-two realization
of the extra weighted slab has been supplied. In its absence, even
minimizing $a^Tx-bx_n^2$ over the bare path, for $b\ge0$, is polynomial:
a concave function has a vertex minimizer, and $x_n^2=L(x)$ at every
vertex, so minimize the linear function $a^Tx-bL(x)$ instead. The slab
creates new vertices and invalidates that linearization.

\subsection{Matrix parameters and nonlinear leaves}\label{s6:parametric}
Interaction rank of a \emph{cost}, treated in Appendix~\ref{app:costs},
must be distinguished from rank of a
parameter perturbing a \emph{constraint matrix}.
Consider rational data and
\[
 \min c^Tx\quad\text{subject to }Ax+\lambda Dx\le b,
 \quad x\in X,\quad\underline\lambda\le\lambda\le\overline\lambda,
\]
where $X$ is a rational polyhedron and the nonempty parameter interval
is finite. If $D=uv^T$, introduce $z=v^Tx$ and $w=\lambda z$.
The exact projected feasible set is the union of the two polyhedra
\begin{align*}
 Ax+uw&\le b,\quad x\in X,\quad z=v^Tx,\quad
 z\ge0,\quad\underline\lambda z\le w\le\overline\lambda z;\\
 Ax+uw&\le b,\quad x\in X,\quad z=v^Tx,\quad
 z\le0,\quad\overline\lambda z\le w\le\underline\lambda z.
\end{align*}
For $z\ne0$, recover $\lambda=w/z$; at $z=0$ both branches force
$w=0$ and any allowed parameter works. Thus two LPs solve joint
feasibility and minimization. If the objective is finite and the set
nonempty, a rational LP solution and that ratio give a rational optimum
of polynomial bit length; unbounded objectives are detected by the LPs.
This auxiliary-variable extension is closely related to the
one-nonzero-column case of Boveroux, Carvalho, Lodi and Louveaux's
preprint \cite[Section~3.3]{s6:boveroux2026}.

Their positive-product construction already uses perturbation rank two:
a fixed-one variable $s$ occurs in the opposing rows for
$\lambda s=z_1$, and a second variable $z_2$ in
$\lambda z_2\le t$. The two nonzero columns have independent supports.
Using the positive-product source in \cref{s3:matsui}, with its explicit
bounds, these rows give ordinary hardness for joint minimization of $t$.
This observation concerns dense base-polytope constraints, not physical
path coefficients. A parameter-dependent right-hand side adds a fixed-one
column before measuring perturbation rank; a direct parameter objective
term and a true max--min objective are outside the two-LP argument.

There is also an arithmetic boundary to replacing the polyhedral blocks
in \cref{s4:fixed-core} by convex nonlinear leaves. Given positive
integers $a_i$ and an integer $B$, let
\[
 P_i=\{y_i:0\le y_i\le\max(1,a_i),\quad y_i^2=a_i\}.
\]
Each leaf is a compact convex singleton in one scalar variable, described by
rational quadratic data, and there is no nonlinear core. The single
aggregate inequality $\sum_i y_i\le B$ asks exactly whether
$\sum_i\sqrt{a_i}\le B$. Hence an unrestricted exact polynomial-bit
extension to these leaves would solve Square-Root Sum in polynomial time.
This is an arithmetic implication, not an $\NP$-hardness theorem.
Independent algebraic leaf values may generate a growing common field,
whereas the polyhedral recovery in \cref{s4:fixed-core} takes place in
the field of the fixed core.

